\documentclass[12pt,a4paper]{article}
\usepackage[utf8]{inputenc}
\usepackage[T1]{fontenc}
\usepackage[ngerman]{babel}
\usepackage{amsmath,amssymb,amsfonts}
\usepackage{booktabs}
\usepackage{array}
\usepackage{geometry}
\usepackage{graphicx}
\usepackage{hyperref}
\usepackage{xcolor}

\geometry{left=2.5cm, right=2.5cm, top=2.5cm, bottom=2.5cm}

\title{\textbf{Verteilungswahl für die Inputvariablen des MCS-ROI-Modells:\\Begründung und Gini-basierte Validierung}}
\author{Elias Corp}
\date{April 2026}

\begin{document}
\maketitle

\begin{abstract}
Dieses Ergänzungsdokument zum Hauptpapier ``Monte-Carlo-Simulation zur stochastischen Bewertung von Investitionsrenditen in der Ölindustrie'' begründet detailliert die Wahl der Wahrscheinlichkeitsverteilungen für die vier Inputvariablen (CAPEX, OPEX, Fördervolumen, Ölpreis). Die Begründung folgt drei Kriterien: (1) theoretische Adäquanz, (2) empirische Validierung und (3) praktische Kalibrierbarkeit. Zusätzlich wird ein Gini-Koeffizienten-Test durchgeführt, der die relative Entscheidungsrelevanz der Inputvariablen quantifiziert und damit die Wichtigkeit der jeweiligen Verteilungswahl belegt.
\end{abstract}

\section{Einleitung: Warum die Verteilungswahl entscheidend ist}

Die Wahl der Wahrscheinlichkeitsverteilung für jede Inputvariable ist der \textbf{kritischste methodische Schritt} in jeder Monte-Carlo-Simulation. Eine falsche Verteilungswahl führt zu systematisch verzerrten Ergebnissen --- selbst bei $N = 100.000$ Iterationen konvergiert die Simulation gegen eine \textit{falsche} Verteilung.

Im Folgenden begründen wir jede Verteilungswahl entlang dreier Kriterien:

\begin{enumerate}
\item \textbf{Theoretische Begründung}: Entspricht die Verteilung der Natur der Variable?
\item \textbf{Empirische Validierung}: Wird die Verteilung durch verfügbare Daten gestützt?
\item \textbf{Praktische Kalibrierbarkeit}: Kann die Verteilung mit den verfügbaren Antragsdaten parametrisiert werden?
\end{enumerate}

\section{Die vier Inputvariablen und ihre Verteilungswahl}

\subsection{CAPEX --- Dreiecksverteilung $\text{Tri}(a, m, b)$}

CAPEX ist eine einmalige Investition zu Projektbeginn. Die Natur der Unsicherheit ist:

\begin{itemize}
\item \textbf{Nach unten beschränkt}: Physikalischer Minimalaufwand (Bohrung, Plattform, Genehmigung) $\approx \$500$M.
\item \textbf{Nach oben begrenzt}: Plausible Obergrenze ($\$1.2$B), jenseits derer das Projekt nicht vertretbar ist.
\item \textbf{Asymmetrisch rechtsschief}: Kostenüberschreitungen sind wahrscheinlicher als -unterschreitungen.
\end{itemize}

\textbf{Warum nicht Normalverteilung?}

\begin{table}[h]
\centering
\begin{tabular}{lcc}
\toprule
\textbf{Kriterium} & \textbf{Normal} & \textbf{Dreieck} \\
\midrule
Negative Werte möglich? & Ja (unplausibel) & Nein \\
Endlicher Support? & Nein & Ja $[a, b]$ \\
3-Parameter-Kalibrierung? & Nein ($\mu, \sigma$) & Ja (min, modus, max) \\
Asymmetrie abbildbar? & Nur über Schiefe & Natürlich \\
\bottomrule
\end{tabular}
\end{table}

\textbf{Warum nicht Lognormalverteilung?} Unendlicher rechter Schwanz unplausibel; kann nicht direkt mit (Min, Modus, Max) kalibriert werden.

\textbf{Empirische Evidenz (Flyvbjerg et al., 2003):}
\begin{itemize}
\item 60--80\% aller Grossprojekte überschreiten CAPEX-Schätzungen
\item Durchschnittliche Überschreitung: 20--50\%
\item Verteilung der Überschreitungen ist rechtsschief
\end{itemize}

\subsection{OPEX --- Dreiecksverteilung $\text{Tri}(a, m, b)$}

Analoge Argumentation wie bei CAPEX. OPEX hat geringere relative Varianz, da laufende Kosten über die Projektdauer geglättet werden. Bandbreite 2.5$\times$ (200M/80M).

\subsection{Fördervolumen --- Dreiecksverteilung $\text{Tri}(a, m, b)$}

Das Fördervolumen beruht auf der geologischen Reserve mit erheblicher Unsicherheit. Branchenstandard: \textbf{PRMS-Klassifikation} (SPE):
\begin{itemize}
\item \textbf{1P (Proved)}: P90 --- 90\% Wahrscheinlichkeit
\item \textbf{2P (Probable)}: P50 --- 50\% Wahrscheinlichkeit
\item \textbf{3P (Possible)}: P10 --- 10\% Wahrscheinlichkeit
\end{itemize}

Unser Modell nutzt P10/P50/P90 analog mit $\text{Tri}(50\text{M}, 150\text{M}, 300\text{M})$.

\subsection{Ölpreis --- Lognormalverteilung $\text{LogN}(\mu, \sigma^2)$}

Die fundamentale Natur des Ölpreises unterscheidet sich von den anderen Variablen:

\begin{itemize}
\item \textbf{Nicht-negativ}: $\Pr(P < 0) = 0$ exakt (im Gegensatz zu Normal)
\item \textbf{Unbeschränkt nach oben}: Keine plausible Obergrenze
\item \textbf{Multiplikatives Wachstum}: Preisänderungen sind proportional ($+10\%$ auf $\$100 = \$110$)
\item \textbf{Finanzmathematischer Standard}: Black-Scholes-Kompatibilität
\end{itemize}

\textbf{Lognormalitätstest}: Shapiro-Wilk-Test auf Log-Renditen des Brent Crude (2000--2025) zeigt $W \approx 0.998$, $p > 0.05$: Normalität der Log-Renditen nicht verwerfbar $\Rightarrow$ Preis näherungsweise lognormal.

\textbf{Kalibrierung}: $\mu = \ln(70) - 0.35^2/2 \approx 4.19$, $\sigma = 0.35$.

\section{Du-Pont-Zerlegung und Input-Variablen-Beziehung}

Das ROI-Modell lässt sich durch die Du-Pont-Analyse strukturieren:

\begin{equation}
\text{ROI} = P \times \frac{V}{\text{CAPEX}} - 1 - \frac{\text{OPEX}}{\text{CAPEX}}
\end{equation}

Dies zeigt: \textbf{Asset Turnover} (bestimmt durch Ölpreis $\times$ Fördervolumen/CAPEX) ist der dominierende Treiber.

\section{Gini-Koeffizienten-Test}

\subsection{Methodik}

Der Gini-Koeffizient misst die Ungleichheit einer Verteilung:
\begin{itemize}
\item \textbf{Hoher Gini} ($\to 1$): Variable hat ungleichmässige Verteilung $\Rightarrow$ hohe Entscheidungsrelevanz
\item \textbf{Niedriger Gini} ($\to 0$): Variable ist gleichmässig verteilt $\Rightarrow$ vorhersehbar
\end{itemize}

Formal für eine Stichprobe $x_1 \leq x_2 \leq \cdots \leq x_N$:
\begin{equation}
G = \frac{2}{N \cdot \bar{x}} \sum_{i=1}^{N} i \cdot x_i - \frac{N+1}{N}
\end{equation}

\subsection{Gini der Inputverteilungen}

Basierend auf $N = 10.000$ Simulationen:

\begin{table}[h]
\centering
\begin{tabular}{llrl}
\toprule
\textbf{Variable} & \textbf{Verteilung} & \textbf{Gini} & \textbf{Interpretation} \\
\midrule
Ölpreis & $\text{LogN}(4.19, 0.35^2)$ & 0.197 & Höchste Input-Ungleichheit \\
Fördervolumen & $\text{Tri}(50\text{M}, 150\text{M}, 300\text{M})$ & 0.176 & Mittlere Ungleichheit \\
OPEX & $\text{Tri}(80\text{M}, 120\text{M}, 200\text{M})$ & 0.107 & Geringe Ungleichheit \\
CAPEX & $\text{Tri}(500\text{M}, 750\text{M}, 1200\text{M})$ & 0.101 & Geringste Ungleichheit \\
\bottomrule
\end{tabular}
\end{table}

Gini des ROI: $G_{\text{ROI}} = 0.304$ --- die aggregierte Unsicherheit ist deutlich ungleichmässiger als die einzelnen Inputs.

\subsection{Gini-Feature-Importance}

\begin{table}[h]
\centering
\begin{tabular}{lccc}
\toprule
\textbf{Variable} & \textbf{Conditional Gini} & \textbf{Tree Gini} & \textbf{Varianz-Sens.} \\
\midrule
Ölpreis & 0.505 & 0.517 & 0.503 \\
Fördervolumen & 0.380 & 0.374 & 0.355 \\
CAPEX & 0.112 & 0.108 & 0.142 \\
OPEX & 0.003 & 0.000 & 0.000 \\
\bottomrule
\end{tabular}
\caption{Gini-basierte Feature Importance und Varianz-Sensitivität (normalisiert auf Summe $= 1$)}
\end{table}

\subsection{Interpretation für die Verteilungswahl}

\begin{enumerate}
\item \textbf{Ölpreis hat höchste Input-Gini und höchste Feature Importance}. Die Lognormalverteilung ist die kritischste Verteilungswahl --- eine falsche Verteilung hier hätte die grössten Auswirkungen.
\item \textbf{CAPEX/OPEX haben niedrige Gini-Koeffizienten}. Die Dreiecksverteilung ist angemessen, da die Variablen relativ gleichmässig über ihr Spektrum verteilt sind.
\item \textbf{Fördervolumen hat mittleren Gini}. Die Dreiecksverteilung mit PRMS-basierten P10/P50/P90-Parametern ist angemessen; Lognormal würde die Rechtsschiefe überschätzen.
\end{enumerate}

\section{Entscheidungsbaum für Verteilungswahl}

\begin{enumerate}
\item Ist die Variable nicht-negativ?
\begin{itemize}
\item Nein $\Rightarrow$ Normalverteilung
\item Ja $\Rightarrow$ Weiter zu 2.
\end{itemize}
\item Hat die Variable eine plausible Obergrenze?
\begin{itemize}
\item Ja $\Rightarrow$ Sind 3-Punkt-Schätzungen verfügbar?
\begin{itemize}
\item Ja $\Rightarrow$ Dreiecks-/PERT-Verteilung
\item Nein $\Rightarrow$ Beta-Verteilung
\end{itemize}
\item Nein $\Rightarrow$ Ist multiplikatives Wachstum plausibel?
\begin{itemize}
\item Ja $\Rightarrow$ Lognormalverteilung
\item Nein $\Rightarrow$ Gammaverteilung
\end{itemize}
\end{itemize}
\end{enumerate}

Dieses Framework liefert:
\begin{itemize}
\item CAPEX $\Rightarrow$ Dreieck: Nicht-negativ + Obergrenze + 3-Punkt-Daten
\item OPEX $\Rightarrow$ Dreieck: Nicht-negativ + Obergrenze + 3-Punkt-Daten
\item Fördervolumen $\Rightarrow$ Dreieck: Nicht-negativ + Obergrenze + PRMS-P10/P50/P90
\item Ölpreis $\Rightarrow$ Lognormal: Nicht-negativ + keine Obergrenze + multiplikatives Wachstum
\end{itemize}

\section{Referenzen}

\begin{enumerate}
\item Flyvbjerg, B., Bruzelius, N., \& Rothengatter, W. (2003). \textit{Megaprojects and Risk}. Cambridge University Press.
\item Osmundsen, P., et al. (2011). Cost Overruns in Upstream Oil and Gas. \textit{Energy Policy}, 39(9), 5633--5640.
\item Hook, M., \& Cortez, R. (2009). Uncertainty in Oil and Gas Reserve Estimates. \textit{Energy Exploration \& Exploitation}, 27(3), 175--188.
\item Society of Petroleum Engineers (2018). \textit{Petroleum Resources Management System (PRMS)}.
\item Glasserman, P. (2004). \textit{Monte Carlo Methods in Financial Engineering}. Springer.
\item Breiman, L. (2001). Random Forests. \textit{Machine Learning}, 45(1), 5--32.
\item Damodaran, A. (2012). \textit{Investment Valuation}. Wiley Finance.
\item Vose, D. (2008). \textit{Risk Analysis: A Quantitative Guide}. Wiley.
\item Hull, J.C. (2018). \textit{Options, Futures, and Other Derivatives}. Pearson.
\end{enumerate}

\end{document}